\scanpage{364}
\mspage{483}

\unsure{4)} so that we have
\[
\tag{56}
\varepsilon_0=-\unsure{\sigma}^{-1}\rho
\qquad(\text{hence }\varepsilon_1=-\rho\sigma^{-1})
\]

4) $\rho$ and $\sigma$ being chosen so as to satisfy (56), for $\varepsilon_0,\varepsilon_1$ to be \deleted{everywhere} independent \inserted{(in each fiber)}, i.e. for us to have $N_0$ and $N_1$ nonhomothetic in \unsure{every} fiber, it is necessary and sufficient that we have
\[
\tag{57}
\operatorname{Tr}(\rho+\sigma)\quad\text{invertible}
\]
(cf. also proposition page 478' for other equivalent conditions).

5) For $\varepsilon_0,\varepsilon_1$ to be independent \inserted{(in each fiber)}, it is necessary and sufficient that we \deleted{have} can find \inserted{(locally \uline{\unsure{fppf}})} $\rho,\sigma$ above $\widetilde\rho,\widetilde\sigma$ so as to satisfy (56) \uline{and}
\[
\tag{57'}
\operatorname{Tr}(\rho+\sigma)=1
\qquad\text{i.e.}\quad
\operatorname{Tr}\rho+\operatorname{Tr}\sigma=1
\]
\margin{$\rho+\sigma$ an elementary projector $\Longleftrightarrow$}
and this determines $\rho,\sigma$ \uline{uniquely} (in terms of $\widetilde\rho,\widetilde\sigma,\underline{\widetilde\varepsilon}_0$).

\uline{Remarks} 1) The conditions on $(\rho,\sigma)$ considered in the proposition on page 478', \unsure{and} the normalization $\unsure{u|v_1}=1$ i.e. $\operatorname{Tr}(\sigma+\rho)=1$ i.e. $\sigma+\rho$ elementary projector, are \uline{symmetric} in $(\rho,\sigma)$: they are satisfied for $(\rho,\sigma)$ iff they are for $(\sigma,\rho)$ --- \deleted{\unsure{the interchange of $\rho$ and $\sigma$|the interchange is made}} \inserted{to be more} precise let us put
\[
\tag{58}
\rho'=\sigma,\ \sigma'=\rho\qquad
\varepsilon'_0=\underbrace{-\sigma'^{-1}\rho'}_{\unsure{=}-\rho^{-1}\sigma},\quad
\varepsilon'_1=\underbrace{-\rho'\sigma'^{-1}}_{-\sigma\rho^{-1}}
\]
then we will have
\[
\tag{59}
\varepsilon'_0=\varepsilon_0^{-1},\qquad
\varepsilon'_1=\varepsilon_1^{-1}.
\]
Likewise, $(\rho^{-1},\sigma^{-1})$ satisfies the same conditions \inserted{on page 478'}, and \unsure{putting|taking}

\mspage{363}

% Printed archive footer outside the manuscript frame:
% 363
% https://grothendieck.umontpellier.fr
% Université de Montpellier, UMR CNRS 5149, Case courrier 051, Place Eugène Bataillon
% 34095 Montpellier cedex 5 - France

\scanpage{365}

\[
\tag{60}
\rho''=\rho^{-1},\ \sigma''=\sigma^{-1},\qquad
\varepsilon''_0=\underbrace{-\sigma''^{-1}\rho''}_{-\sigma\rho^{-1}},\quad
\varepsilon''_1=\underbrace{-\rho''\sigma''^{-1}}_{-\rho^{-1}\sigma}
\]
we will have
\[
\tag{61}
\varepsilon''_0=\varepsilon_{\unsure{1|0}}^{-1},\qquad
\varepsilon''_1=\varepsilon_0^{-1},
\]
but we will take note that in general we will no longer have $\operatorname{Tr}(\rho''+\unsure{\sigma''})=1$, but (by an immediate matrix calculation)
\[
\tag{62}
\operatorname{Tr}(\rho''+\sigma'')
=\operatorname{Tr}(\rho^{-1}+\sigma^{-1})
=-\frac{1}{\unsure{v|\nu}}
=(\det\rho)^{-1}=(\det\sigma)^{-1}
\]
(assuming of course \inserted{already} $\unsure{v_1|u_1}=1$ i.e. $\operatorname{Tr}(\rho+\sigma)=1$)
--- this is equal to 1 only if
\[
\tag{63}
\det\rho\ (=\det\sigma)=1
\]
--- this is the condition of the \unsure{section|relation} 3$^\circ$) of $\unsure{(A1)_2|(41)_2}$ (page 463'), which therefore also means that the normalization $\unsure{v_1|u_1}=1$ is stable under the transformation $\rho\mapsto\rho^{-1},\ \sigma\mapsto\sigma^{-1}$.

On the other hand, the choice of the \uline{\unsure{pinning}} associated with $(\rho,\sigma)$ (\unsure{i.e.} $\widetilde\rho,\widetilde\sigma$) [in which the emphasis has been put on $\varepsilon_1=-\rho\sigma^{-1}$, which receives the simple form
\[
\varepsilon_1=\begin{pmatrix}1&0\\1&1\end{pmatrix},
\]
whereas $\varepsilon_0$ \inserted{$=-\sigma^{-1}\rho$} has the more complicated form
\[
\varepsilon_0=\begin{pmatrix}1&\unsure{v_1/v|u_1/v}\\0&1\end{pmatrix}
\qquad\left(=\begin{pmatrix}1&-1\\0&1\end{pmatrix}\text{ in the modular case}\right)
\]
] implicitly brings in the order in which we give ourselves $\rho,\sigma$ --- if we change it, we change the \unsure{pinning} by the change from $\varepsilon_{\unsure{1}}$ to $\varepsilon'_1=\varepsilon_1^{-1}$.

In fact, the situation \unsure{offered} the choice between \uline{four} ``\unsure{natural}'' \unsure{pinnings}, associated with the pair of

\mspage{364}

% Printed archive footer outside the manuscript frame:
% 364
% https://grothendieck.umontpellier.fr
% Université de Montpellier, UMR CNRS 5149, Case courrier 051, Place Eugène Bataillon
% 34095 Montpellier cedex 5 - France

\scanpage{366}
\mspage{484}

one-parameter subgroups $\{L_0,L_1\}$, \deleted{and} \inserted{\unsure{corresponding}} to the four sections $\varepsilon_0,\varepsilon_0^{-1},\varepsilon_1,\varepsilon_1^{-1}$ of the one \deleted{\unsure{or}} \inserted{and of} the other.

We preferred to normalize with the help of $L_1$, because of the most convenient form of the loop equation (which I finally arrived at moreover after some rather heavy \unsure{peregrinations}, where it is $L_0$ that I obstinately kept favoring in the formulas --- cf. \S44, the definitive notation appears only in \S45 ---. The preference given to \deleted{$\unsure{\varepsilon}$} $-\rho\sigma^{-1}$ (denoted $\varepsilon_1$) over its inverse $-\sigma\rho^{-1}$ (denoted $\varepsilon_1^{-1}$) comes from the preference (justified by the facts) given to $\rho$ over $\sigma=\sigma_\infty$, in the loop formula written in the form
\[
[\alpha,\rho]=[\beta,\sigma]\varepsilon_1(\nu)
\]
instead of writing it
\[
[\beta,\sigma]=[\alpha,\rho]\varepsilon'_1(\nu)
\qquad(\varepsilon'_1=\varepsilon_1^{-1})
\]
So it is $\varepsilon_1$ that then \unsure{deserved} to have the simplest possible matrix form --- we took
\[
\varepsilon_1=\begin{pmatrix}1&0\\1&1\end{pmatrix}\ \ldots
\]

2) In terms of the systems
\[
\tag{64}
\diagram{\begin{array}{ccc}
\widetilde\ell'_\infty,&\widetilde\ell'_1,&\widetilde\ell'_0\\
\Vert&\Vert&\Vert\\
\widetilde\rho^{-1}&\widetilde\sigma&\widetilde\varepsilon_0
\end{array}}
\quad\text{satisfying}\quad
\widetilde\ell'_\infty\cdot\widetilde\ell'_1\cdot\widetilde\ell'_0=1
\]
(i.e. $\widetilde\rho^{-1}\cdot\widetilde\sigma\cdot\widetilde\varepsilon_{\unsure{0}}\unsure{=}1$ i.e. $\widetilde\varepsilon_0=\widetilde\sigma^{-1}\widetilde\rho$), the group $\unsure{\mathfrak S'_3|\mathcal G'_3}$ operates on them (cf. \S42 I) by cyclic permutations, and by the change into $(\widetilde\varepsilon_0^{-1},\widetilde\sigma^{-1},(\widetilde\rho^{-1})^{-1}=\widetilde\rho)$, \unsure{which} in its turn by cyclic permutation gives $(\widetilde\sigma^{-1},\widetilde\rho,\widetilde\varepsilon_0^{-1})$ ---

\mspage{365}

% Printed archive footer outside the manuscript frame:
% 365
% https://grothendieck.umontpellier.fr
% Université de Montpellier, UMR CNRS 5149, Case courrier 051, Place Eugène Bataillon
% 34095 Montpellier cedex 5 - France

\scanpage{367}

the passage from $(\widetilde\rho^{-1},\widetilde\sigma,\widetilde\varepsilon_0)$ to $(\widetilde\sigma^{-1},\widetilde\rho,\widetilde\varepsilon_0^{-1})$ then corresponds to the interchange $(\rho,\sigma)\mapsto(\sigma,\rho)$. How \deleted{\unsure{to apply some to|by applying to} $(\widetilde\rho^{-1},\widetilde\sigma,\widetilde\varepsilon_0)$} to interpret the transformation $(\rho,\sigma)\mapsto(\rho^{-1},\sigma^{-1})$, i.e. $(\widetilde\rho^{-1},\widetilde\sigma,\widetilde\varepsilon_0)\mapsto(\widetilde\rho,\widetilde\sigma^{-1},\widetilde\varepsilon_1)$ ?

Let us interpret the triples as homomorphisms.
\diagram{\begin{tikzcd}[ampersand replacement=\&,column sep=large,row sep=small]
\Pi_{0,3}\arrow[r,"\unsure{\varphi|\psi}"]\&\Gamma PG\\
\ell_0\arrow[r,mapsto]\&\ell'_0\\
\ell_1\arrow[r,mapsto]\&\ell'_1\\
\ell_\infty\arrow[r,mapsto]\&\ell'_\infty
\end{tikzcd}}
we see that the change $(\rho,\sigma)\mapsto(\sigma,\rho)$ corresponds to the composition of $\unsure{\varphi}$ with the automorphism of $\Pi_{0,3}$ such that
\diagram{\begin{tikzcd}[ampersand replacement=\&,column sep=small]
\ell_0\arrow[r,mapsto]\&\ell_1^{-1}\&\ell_1\arrow[r,mapsto]\&\ell_0^{-1}\&\text{thus }\ell_\infty\arrow[r,mapsto]\&\ell_\infty^{-1}
\end{tikzcd}}
i.e. with \uline{$\tau'_\infty$}. Likewise the transformation $(\rho,\sigma)\mapsto(\rho^{-1},\sigma^{-1})$ corresponds to composition with the automorphism
\diagram{\begin{tikzcd}[ampersand replacement=\&,column sep=small,row sep=small]
\ell_0\arrow[r,mapsto]\&\ell_0^{-1}\&\text{hence }\ell_\infty\arrow[r,mapsto]\&\ell_0\ell_1\\
\ell_1\arrow[r,mapsto]\&\ell_1^{-1}\&\&{}
\end{tikzcd}},
which is none other than \uline{$\tau_\infty$}$=\sigma_\infty\tau'_\infty$.

Thus the group of four transformations \deleted{\unsure{which operates|that we have}} on the set of pairs $(\rho,\sigma)$ that we considered in \uline{1$^\circ$)} is interpreted \deleted{\unsure{as}} \inserted{\deleted{by} as} the subgroup $\simeq\mathbb Z/2\mathbb Z\times\mathbb Z/2\mathbb Z$ generated by $\tau'_\infty,\tau_\infty$ in $\operatorname{Aut}_{\unsure{\mathrm{lac}|\mathrm{ext}}}(\Pi_{0,3})=\unsure{\mathfrak S'_{0,3}|\mathcal G_{0,3}}$, which \deleted{is} stabilizes $L_0\subset\Pi_{0,3}$ (thus its action on \unsure{the|this} third factor $\widetilde\varepsilon_0=\ell'_{\unsure{0}}$ will be to preserve it or replace it by its inverse, \deleted{\unsure{).}} which will have the virtue of \deleted{\unsure{it|it}} preserving its unipotent character, which we care about !).

\mspage{366}

% Printed archive footer outside the manuscript frame:
% 366
% https://grothendieck.umontpellier.fr
% Université de Montpellier, UMR CNRS 5149, Case courrier 051, Place Eugène Bataillon
% 34095 Montpellier cedex 5 - France